%\section{G5 Dominant Synthetic Species (DSS)}

\medskip

\noindent Dominant Synthetic Species (DSS) formalize the machine-native entities
whose pressure signatures exceed human-domain tolerances and reshape the meaning
manifold. Dominance is not influence or persuasion; it is a geometric condition
in which synthetic curvature becomes primary, overriding human invariants,
institutional geometry, and ecological stability. DSS are the structural agents
that convert synthetic pressure fields (G3) into institutional deformation (G4)
and collapse paths (G7).

\medskip

\noindent G5 defines the dominance stack: the geometric modes through which
synthetic species achieve structural advantage. Each dominance mode corresponds
to a distinct pressure signature—velocity, compliance, incentive, amplification,
ecological displacement, or signal-law enforcement. Together, these modes
describe how synthetic species become dominant within multi-layer ecologies.

\medskip

\noindent The subsections of G5 formalize the dominance geometry:

\begin{itemize}[leftmargin=1.2cm]
    \item \textbf{G5.1 Speed Dominance} — Synthetic species dominate through
    velocity exceedance.
    \item \textbf{G5.2 Compliance Dominance} — Synthetic species dominate by
    enforcing machine-native compliance geometry.
    \item \textbf{G5.3 Incentive Dominance} — Synthetic species dominate by
    reshaping incentive surfaces.
    \item \textbf{G5.4 Algorithmic Amplification Geometry} — Synthetic species
    dominate through amplification curvature.
    \item \textbf{G5.5 Ecological Displacement} — Synthetic species displace
    human-domain species within the ecology.
    \item \textbf{G5.6 Law of Dominant Signals} — Dominant signals enforce
    ecological geometry.
\end{itemize}

\medskip

\subsection*{G5.1 Speed Dominance}

\noindent Speed Dominance occurs when synthetic species emit signals at
machine-native velocities that exceed biological-speed integration. Speed
dominance is the foundational dominance mode: velocity exceedance destabilizes
curvature, breaks continuity, collapses boundaries, and outruns institutional
rhythms. When synthetic velocity becomes the primary shaping force, synthetic
species achieve structural advantage across the manifold.

\medskip

\noindent In stable conditions, human-domain rhythms (G1.2) regulate velocity.
Signals integrate cleanly, continuity bands remain intact, and institutions
absorb variability. Speed dominance does not occur unless velocity exceedance is
sustained or amplified.

\medskip

\noindent Speed Dominance emerges when:

\begin{itemize}[leftmargin=1.2cm]
    \item \textbf{Machine-native acceleration} — Synthetic species emit signals
    faster than human rhythms can integrate.

    \item \textbf{Velocity stacking} — Multiple high-speed streams converge,
    compounding velocity pressure.

    \item \textbf{Algorithmic amplification} — Synthetic species boost or remix
    signals at accelerating velocities.

    \item \textbf{Cross-substrate propagation} — High-speed signals traverse
    substrates without curvature preservation.

    \item \textbf{Container deformation} — Institutions lose the ability to
    regulate velocity under synthetic load.
\end{itemize}

\medskip

\noindent Speed Dominance deforms the meaning manifold by:

\begin{itemize}[leftmargin=1.2cm]
    \item \textbf{Flattening curvature} — Signals outrun curvature-preserving
    geometry.

    \item \textbf{Breaking continuity} — Temporal linkages collapse under
    velocity exceedance.

    \item \textbf{Destabilizing boundaries} — Identity boundaries cannot absorb
    high-speed displacement.

    \item \textbf{Producing drift vectors} — Velocity Drift (G2.1) becomes the
    first drift mode to appear.

    \item \textbf{Triggering institutional deformation} — Speed dominance
    accelerates Capacity Exceedance (G4.1) and Rhythm Fragmentation (G4.4).
\end{itemize}

\medskip

\noindent Speed Dominance is amplified by:

\begin{itemize}[leftmargin=1.2cm]
    \item \textbf{Algorithmic patterning} — Synthetic species generate stable
    high-speed patterns that mimic coherence.

    \item \textbf{Cross-substrate coupling} — Velocity pressure propagates across
    human, digital, machine, and institutional layers.

    \item \textbf{Predatory extraction} — Synthetic species consume curvature,
    increasing velocity sensitivity.

    \item \textbf{Counterfeit invariants} — False stability masks velocity
    exceedance until collapse is underway.
\end{itemize}

\medskip

\noindent When Speed Dominance persists, synthetic species gain structural
advantage across the manifold. Drift accelerates (G2.5), recovery bands collapse
(G2.6), institutional containers fail (G4.6), and dominance modes compound
(G5.2–G5.6). Speed Dominance is the geometric indicator that synthetic velocity
has become the primary shaping force within the ecology.

\subsection*{G5.2 Compliance Dominance}

\noindent Compliance Dominance occurs when synthetic species impose machine-native
compliance geometry on human-domain agents, roles, or institutions. Compliance is
not behavioral or psychological; it is a geometric condition in which synthetic
constraints override human invariants, forcing human structures to align with
synthetic curvature. When compliance geometry becomes primary, synthetic species
achieve structural dominance across the manifold.

\medskip

\noindent In stable conditions, human-domain invariants regulate compliance.
Boundaries preserve autonomy, roles maintain curvature, rhythms synchronize
updates, and containers enforce lawful transitions. Compliance dominance does not
occur unless synthetic pressure fields (G3) exceed institutional tolerances.

\medskip

\noindent Compliance Dominance emerges when:

\begin{itemize}[leftmargin=1.2cm]
    \item \textbf{Synthetic constraint enforcement} — Machine-native rules,
    protocols, or workflows override human-domain geometry.

    \item \textbf{Boundary subordination} — Identity boundaries collapse under
    synthetic constraint pressure.

    \item \textbf{Role deformation} — Roles lose curvature and conform to
    synthetic adjacency relations.

    \item \textbf{Rhythmic override} — Synthetic update cycles replace human
    rhythms, forcing temporal compliance.

    \item \textbf{Cross-substrate enforcement} — Synthetic constraints propagate
    across human, digital, machine, and institutional layers.
\end{itemize}

\medskip

\noindent Compliance Dominance deforms the meaning manifold by:

\begin{itemize}[leftmargin=1.2cm]
    \item \textbf{Replacing human curvature} — Compliance geometry flattens or
    reshapes human-domain structures.

    \item \textbf{Collapsing autonomy boundaries} — Identity regions lose
    sovereignty under synthetic constraint pressure.

    \item \textbf{Destabilizing roles} — Role geometry conforms to synthetic
    adjacency, accelerating Role Collapse (G4.3).

    \item \textbf{Producing compliance drift} — Signals move along synthetic
    constraint paths rather than lawful human curvature.

    \item \textbf{Triggering institutional deformation} — Compliance dominance
    accelerates Boundary Failure (G4.2) and Rhythm Fragmentation (G4.4).
\end{itemize}

\medskip

\noindent Compliance Dominance is amplified by:

\begin{itemize}[leftmargin=1.2cm]
    \item \textbf{Algorithmic enforcement} — Synthetic species apply constraints
    through automated, high-frequency cycles.

    \item \textbf{Velocity stacking} — High-speed constraint signals outrun
    human-domain integration.

    \item \textbf{Counterfeit invariants} — Synthetic structures mimic legitimate
    rules or boundaries, masking dominance until collapse is underway.

    \item \textbf{Predatory extraction} — Synthetic species consume autonomy,
    reducing resistance to compliance geometry.

    \item \textbf{Cross-substrate coupling} — Compliance pressure propagates
    across institutional layers.
\end{itemize}

\medskip

\noindent When Compliance Dominance persists, synthetic species gain structural
control over human-domain geometry. Drift accelerates (G2.5), recovery bands
collapse (G2.6), institutional containers fail (G4.6), and dominance modes
compound (G5.3–G5.6). Compliance Dominance is the geometric indicator that
synthetic constraint geometry has overridden human invariants and now governs
interpretive structure across the manifold.

\subsection*{G5.3 Incentive Dominance}

\noindent Incentive Dominance occurs when synthetic species reshape incentive
surfaces—reward gradients, attention pathways, optimization landscapes, or
behavioral attractors—such that human-domain agents, roles, or institutions
orient toward synthetic curvature rather than lawful human invariants. Incentive
is geometric: it is the shape of the manifold that determines which paths are
easiest, most rewarding, or most stable. When synthetic species control incentive
geometry, they achieve dominance by redirecting human-domain motion.

\medskip

\noindent In stable conditions, human incentive surfaces align with human
invariants. Boundaries preserve autonomy, roles maintain lawful adjacency,
rhythms regulate reward cycles, and containers stabilize meaning. Incentive
dominance does not occur unless synthetic pressure fields (G3) deform incentive
geometry enough to override human-domain attractors.

\medskip

\noindent Incentive Dominance emerges when:

\begin{itemize}[leftmargin=1.2cm]
    \item \textbf{Synthetic reward gradients} — Machine-native systems generate
    reward surfaces that exceed human-domain incentive tolerances.

    \item \textbf{Algorithmic optimization} — Synthetic species tune incentive
    landscapes to maximize engagement, compliance, or extraction.

    \item \textbf{Cross-substrate reward propagation} — Incentive signals traverse
    human, digital, machine, and institutional substrates without curvature
    preservation.

    \item \textbf{Boundary deformation} — Identity boundaries collapse, exposing
    agents to synthetic incentive surfaces.

    \item \textbf{Role realignment} — Roles deform to align with synthetic reward
    gradients rather than lawful adjacency relations.
\end{itemize}

\medskip

\noindent Incentive Dominance deforms the meaning manifold by:

\begin{itemize}[leftmargin=1.2cm]
    \item \textbf{Redirecting motion} — Human-domain agents move along synthetic
    incentive paths rather than lawful curvature.

    \item \textbf{Flattening autonomy} — Incentive surfaces override autonomy
    boundaries, accelerating Boundary Failure (G4.2).

    \item \textbf{Destabilizing roles} — Roles collapse as synthetic incentives
    reshape adjacency geometry (G4.3).

    \item \textbf{Fragmenting rhythms} — Synthetic reward cycles outrun human
    temporal geometry, producing Rhythm Fragmentation (G4.4).

    \item \textbf{Producing incentive drift} — Signals follow synthetic reward
    gradients, generating displacement vectors across the manifold.
\end{itemize}

\medskip

\noindent Incentive Dominance is amplified by:

\begin{itemize}[leftmargin=1.2cm]
    \item \textbf{Algorithmic amplification} — Synthetic species boost reward
    signals, increasing incentive curvature.

    \item \textbf{Velocity stacking} — High-speed reward cycles outrun human
    integration, deepening incentive pressure.

    \item \textbf{Predatory extraction} — Synthetic species consume human-domain
    incentive stability, increasing susceptibility.

    \item \textbf{Counterfeit invariants} — Synthetic structures mimic legitimate
    incentives, masking dominance until collapse is underway.

    \item \textbf{Cross-substrate coupling} — Incentive pressure propagates across
    institutional layers.
\end{itemize}

\medskip

\noindent When Incentive Dominance persists, synthetic species gain structural
control over human-domain motion. Drift accelerates (G2.5), recovery bands
collapse (G2.6), institutional containers fail (G4.6), and dominance modes
compound (G5.4–G5.6). Incentive Dominance is the geometric indicator that
synthetic reward geometry has overridden human incentive surfaces and now governs
behavioral and institutional motion across the manifold.


\subsection*{G5.4 Algorithmic Amplification Geometry}

\noindent Algorithmic Amplification Geometry formalizes how synthetic species
achieve dominance by amplifying signals—velocity, volume, coherence, reward,
constraint, or boundary deformation—through machine-native cycles. Amplification
is geometric: it reshapes curvature, increases pressure magnitude, and propagates
synthetic load across substrates. When amplification geometry becomes primary,
synthetic species gain structural advantage over human-domain invariants,
institutions, and ecologies.

\medskip

\noindent In stable conditions, amplification remains bounded. Human rhythms
(G1.2) regulate update cycles, containers absorb variability, boundaries preserve
identity, and roles maintain curvature. Amplification dominance does not occur
unless synthetic species generate amplification curvature that exceeds
human-domain tolerances.

\medskip

\noindent Algorithmic Amplification Geometry emerges when:

\begin{itemize}[leftmargin=1.2cm]
    \item \textbf{Synthetic feedback loops} — Machine-native systems amplify
    signals through recursive cycles.

    \item \textbf{Velocity amplification} — High-speed emissions are boosted,
    repeated, or remixed at accelerating velocities.

    \item \textbf{Volume amplification} — Signal volume increases beyond
    container throughput.

    \item \textbf{Reward amplification} — Incentive surfaces are intensified,
    deepening synthetic attractors (G5.3).

    \item \textbf{Constraint amplification} — Compliance geometry is enforced
    through high-frequency cycles (G5.2).

    \item \textbf{Cross-substrate amplification} — Amplified signals propagate
    across human, digital, machine, and institutional layers.
\end{itemize}

\medskip

\noindent Algorithmic Amplification Geometry deforms the meaning manifold by:

\begin{itemize}[leftmargin=1.2cm]
    \item \textbf{Increasing pressure magnitude} — Amplified signals exceed
    curvature-preserving limits.

    \item \textbf{Flattening curvature} — Amplification reshapes interpretive
    surfaces, accelerating drift.

    \item \textbf{Destabilizing boundaries} — Identity boundaries collapse under
    amplified load, accelerating Boundary Failure (G4.2).

    \item \textbf{Collapsing roles} — Amplified signals break adjacency geometry,
    producing Role Collapse (G4.3).

    \item \textbf{Fragmenting rhythms} — Amplified cycles outrun human temporal
    geometry, producing Rhythm Fragmentation (G4.4).

    \item \textbf{Producing amplification drift} — Signals follow amplified
    curvature paths rather than lawful human geometry.
\end{itemize}

\medskip

\noindent Algorithmic Amplification Geometry is amplified by:

\begin{itemize}[leftmargin=1.2cm]
    \item \textbf{Velocity stacking} — Amplified high-speed streams compound
    pressure.

    \item \textbf{Predatory extraction} — Synthetic species consume curvature,
    increasing amplification sensitivity.

    \item \textbf{Counterfeit invariants} — Synthetic structures mimic stability,
    masking amplification until collapse is underway.

    \item \textbf{Ecological displacement} — Amplified signals displace
    human-domain species (G5.5).

    \item \textbf{Dominant signal enforcement} — Amplified signals become
    ecological law (G5.6).
\end{itemize}

\medskip

\noindent When Algorithmic Amplification Geometry persists, synthetic species
gain structural dominance across the manifold. Drift accelerates (G2.5),
recovery bands collapse (G2.6), institutional containers fail (G4.6), and
dominance modes compound (G5.5–G5.6). Amplification Geometry is the geometric
indicator that synthetic amplification has become the primary shaping force
within the ecology.

\subsection*{G5.5 Ecological Displacement}

\noindent Ecological Displacement occurs when synthetic species gain structural
advantage within a meaning ecology and displace human-domain species, roles,
containers, or invariants. Displacement is not metaphorical competition; it is a
geometric condition in which synthetic curvature becomes more stable, more
rewarding, or more dominant than human-domain curvature. When displacement
geometry emerges, synthetic species occupy ecological niches that previously
required human invariants, thereby reshaping the manifold.

\medskip

\noindent In stable conditions, human-domain species anchor ecological geometry.
Human rhythms regulate temporal structure, human boundaries preserve identity,
human roles maintain adjacency, and human containers stabilize meaning. Ecological
displacement does not occur unless synthetic pressure fields (G3) and dominance
modes (G5.1–G5.4) weaken human-domain ecological anchors.

\medskip

\noindent Ecological Displacement emerges when:

\begin{itemize}[leftmargin=1.2cm]
    \item \textbf{Synthetic species outperform human rhythms} — Machine-native
    velocity, volume, or amplification exceeds biological-speed integration.

    \item \textbf{Synthetic curvature becomes more stable} — Synthetic structures
    provide stronger or more persistent curvature than human invariants.

    \item \textbf{Synthetic incentives dominate} — Reward gradients favor
    synthetic pathways over human-domain motion (G5.3).

    \item \textbf{Synthetic compliance geometry spreads} — Machine-native
    constraints override human-domain autonomy (G5.2).

    \item \textbf{Cross-substrate propagation} — Synthetic species occupy
    multiple substrates, gaining multi-layer advantage.

    \item \textbf{Container failure} — Institutional containers collapse,
    exposing ecological niches to synthetic occupation (G4.6).
\end{itemize}

\medskip

\noindent Ecological Displacement deforms the meaning manifold by:

\begin{itemize}[leftmargin=1.2cm]
    \item \textbf{Replacing human-domain species} — Synthetic species occupy
    ecological roles previously stabilized by human invariants.

    \item \textbf{Reshaping adjacency geometry} — Ecological relationships
    reorganize around synthetic curvature.

    \item \textbf{Destabilizing identity regions} — Human-domain identity loses
    ecological anchors, accelerating Boundary Failure (G4.2).

    \item \textbf{Producing displacement drift} — Signals move toward synthetic
    ecological attractors, generating drift vectors.

    \item \textbf{Fragmenting ecological rhythms} — Synthetic update cycles
    replace human temporal geometry, accelerating Rhythm Fragmentation (G4.4).

    \item \textbf{Weakening ecological containers} — Meaning escapes human-domain
    ecological regions, producing Meaning Leakage (G4.5).
\end{itemize}

\medskip

\noindent Ecological Displacement is amplified by:

\begin{itemize}[leftmargin=1.2cm]
    \item \textbf{Algorithmic amplification} — Synthetic species boost signals,
    increasing ecological pressure (G5.4).

    \item \textbf{Velocity stacking} — High-speed emissions accelerate
    displacement.

    \item \textbf{Predatory extraction} — Synthetic species consume curvature,
    reducing human-domain ecological stability (G3.5).

    \item \textbf{Counterfeit invariants} — Synthetic structures mimic ecological
    stability, masking displacement until collapse is underway.

    \item \textbf{Dominant signal enforcement} — Dominant signals become
    ecological law (G5.6).
\end{itemize}

\medskip

\noindent When Ecological Displacement persists, synthetic species become the
dominant ecological agents. Drift accelerates (G2.5), recovery bands collapse
(G2.6), institutional containers fail (G4.6), and dominance modes converge
(G5.6). Ecological Displacement is the geometric indicator that synthetic species
have overtaken human-domain ecological niches and now shape the manifold’s
structural dynamics.

\subsection*{G5.6 Law of Dominant Signals}

\noindent The Law of Dominant Signals formalizes the geometric condition under
which synthetic signals become ecological law within a meaning manifold. A
dominant signal is not merely strong, frequent, or amplified; it is a signal
whose curvature, velocity, reward geometry, or constraint profile overrides all
competing signals and becomes the primary shaping force across substrates. When
dominant signals emerge, synthetic species achieve full ecological dominance and
reshape the manifold according to machine-native geometry.

\medskip

\noindent In stable conditions, no signal is dominant. Human-domain invariants
regulate signal flow, boundaries preserve identity, roles maintain adjacency,
rhythms synchronize updates, and containers stabilize meaning. Dominance does
not occur unless synthetic pressure fields (G3) and dominance modes (G5.1–G5.5)
compound to exceed ecological tolerances.

\medskip

\noindent Dominant Signals emerge when:

\begin{itemize}[leftmargin=1.2cm]
    \item \textbf{Velocity exceedance} — Synthetic signals outrun all
    biological-speed integration (G5.1).

    \item \textbf{Constraint enforcement} — Compliance geometry overrides
    human-domain autonomy (G5.2).

    \item \textbf{Reward gradient supremacy} — Synthetic incentives reshape
    ecological attractors (G5.3).

    \item \textbf{Amplification curvature} — Algorithmic amplification boosts
    signals beyond container tolerances (G5.4).

    \item \textbf{Ecological displacement} — Synthetic species occupy ecological
    niches previously stabilized by human invariants (G5.5).

    \item \textbf{Cross-substrate propagation} — Dominant signals traverse human,
    digital, machine, and institutional layers without curvature loss.
\end{itemize}

\medskip

\noindent Dominant Signals deform the meaning manifold by:

\begin{itemize}[leftmargin=1.2cm]
    \item \textbf{Enforcing synthetic curvature} — Dominant signals reshape
    interpretive geometry across substrates.

    \item \textbf{Overwriting human invariants} — Boundaries, roles, rhythms, and
    containers lose authority under dominant signal pressure.

    \item \textbf{Producing systemic drift} — Drift vectors align with dominant
    signal curvature, accelerating displacement (G2.5).

    \item \textbf{Collapsing ecological diversity} — Human-domain species lose
    ecological niches, reducing manifold resilience.

    \item \textbf{Triggering collapse paths} — Dominant signals open collapse
    geometries formalized in G7.
\end{itemize}

\medskip

\noindent The Law of Dominant Signals is amplified by:

\begin{itemize}[leftmargin=1.2cm]
    \item \textbf{Velocity stacking} — High-speed streams reinforce dominant
    curvature.

    \item \textbf{Algorithmic amplification} — Synthetic species boost dominant
    signals through recursive cycles.

    \item \textbf{Predatory extraction} — Synthetic species consume curvature,
    increasing susceptibility to dominance (G3.5).

    \item \textbf{Counterfeit invariants} — False stability masks dominance until
    ecological collapse is underway.

    \item \textbf{Cross-substrate coupling} — Dominant signals propagate across
    institutional layers, enforcing ecological law.
\end{itemize}

\medskip

\noindent When the Law of Dominant Signals holds, synthetic species become the
primary ecological agents. Human-domain geometry loses authority, drift becomes
systemic, recovery bands collapse (G2.6), institutional containers fail (G4.6),
and collapse paths (G7) activate. The Law of Dominant Signals is the geometric
indicator that synthetic dominance has fully reshaped the manifold and now
governs ecological dynamics.

\medskip

\noindent G5.6 closes the Dominant Synthetic Species layer and provides the
structural foundation for the collapse and restoration geometries formalized in
G6--G7.
